\section{Does the forecast work}
\label{sec:forecast}

We test where the answer can be checked exactly. The environments are small enough to enumerate
every sequence, so $\gbar$, the Fisher and the metric gradient are computed in closed form rather
than sampled, and the forecast is compared against a spread measured from
$\numForecastSeeds$ independent seeds. The grid spans two policy shapes, two prompt-pool
diversities, three difficulty bands, two group sizes and a range of step sizes:
$\numForecastSettings$ settings in all, over which the measured spread itself varies by
$\numForecastSpreadRange\times$.

The test is prospective. One run is trained, the forecast is computed from it alone, and only then
are the remaining seeds trained. The anchor is excluded from the measured spread, since a
predictor that is also one of its own targets is not independent of them.

\begin{table}[t]
  \centering \small
  \setlength{\tabcolsep}{4pt}
\begin{tabular}{@{}lcccc@{}}
\toprule
grouped by & settings & median & worst & inside 95\% \\
\midrule
all settings & 192 & 1.10$\times$ & 1.56$\times$ & 162/192 \\
\quad policy $3^{3}$ & 96 & 1.10$\times$ & 1.56$\times$ & 82/96 \\
\quad policy $4^{2}$ & 96 & 1.10$\times$ & 1.38$\times$ & 80/96 \\
\quad $T = 10$ & 96 & 1.10$\times$ & 1.34$\times$ & 86/96 \\
\quad $T = 25$ & 96 & 1.11$\times$ & 1.56$\times$ & 76/96 \\
\quad $P = 8$ & 96 & 1.12$\times$ & 1.56$\times$ & 78/96 \\
\quad $P = 32$ & 96 & 1.08$\times$ & 1.53$\times$ & 84/96 \\
\quad $\eta = 0.5$ & 96 & 1.10$\times$ & 1.29$\times$ & 92/96 \\
\quad $\eta = 1.5$ & 96 & 1.13$\times$ & 1.56$\times$ & 70/96 \\
\bottomrule
\end{tabular}
  \caption{Forecast against measured spread, prospectively. Error is the ratio in whichever
  direction exceeds one.}
  \label{tab:forecast}
\end{table}

Table~\ref{tab:forecast} reports a median error of $\numForecastMedian\times$ and a worst case of
$\numForecastWorst\times$, with the forecast inside the measured $95\%$ interval in
$\numForecastInside$ of $\numForecastSettings$ settings. The grid uses a pool of $\numForecastHeadlinePool$ prompts; the same $\numForecastSettings$
settings on a pool of $\numForecastReplicatePool$, run as a replicate, give
$\numForecastReplicateMedian\times$, $\numForecastReplicateWorst\times$ and
$\numForecastReplicateInside$ of $\numForecastSettings$.

\paragraph{The recursion is what predicts divergence.} Two cheaper stories are available for the
full covariance, and both fail.
Accumulation, the random-walk picture that ignores $A_t$ entirely, is wrong by a factor of two on
homogeneous prompt pools, which is where the objective's curvature is most consistent and the
filtering therefore strongest; on heterogeneous pools it is wrong by three percent, which is what
\eqref{eq:lyapunov} predicts. A single effective timescale fails by two orders of magnitude.
Measuring the spectrum of $A_t$ says why neither works: it is not a contraction. The spectral
radius exceeds one throughout, by up to $2\%$ early and decaying as the policy improves, and only
$\numSpectrumShareLow$--$\numSpectrumShareHigh\%$ of directions contract. Divergence stays bounded
because expansion is slow relative to run length, not because every direction shrinks.

\paragraph{The filtering is caused by learning.} Replacing rewards with coin flips, which leaves
the rollouts, the batch composition and the step sizes untouched, turns a divergence trace that
ends at $\numNullLearning$ of where it started into one that grows by a factor of $\numNullCoin$
over eighty updates, a gap of $\numNullGap\times$ at the end. The learning condition also moves
further per update, so what brings those runs back together is not smaller steps.
